\section{Behavior Cloning y DAgger en el problema del acantilado}

\subsection*{a) Ecuaciones de Bellman para el costo descontado}

Denotemos por $S_\pi$ y $D_\pi$ los costos descontados esperados al
comenzar en el camino seguro y en la zanja, respectivamente. Así,
$J(\pi)=S_\pi$. Desde el camino seguro, ambas políticas satisfacen
\[
    S_\pi=(1-\epsilon)\gamma S_\pi
          +\epsilon(1+\gamma D_\pi).
\]
El costo de caída $1$ se paga en el paso que lleva a la zanja.

Con Behavior Cloning, la zanja es absorbente y cada transición desde
ella cuesta $1$. Por tanto,
\begin{align*}
    D_{\mathrm{BC}} &= 1+\gamma D_{\mathrm{BC}}
        =\frac{1}{1-\gamma}, \\
    \bigl(1-\gamma(1-\epsilon)\bigr)S_{\mathrm{BC}}
        &=\epsilon(1+\gamma D_{\mathrm{BC}})
        =\frac{\epsilon}{1-\gamma}, \\
    J(\pi_{\mathrm{BC}})&=
        \boxed{\frac{\epsilon}{(1-\gamma)(1-\gamma+\gamma\epsilon)}}.
\end{align*}

Con DAgger, la recuperación lleva al camino seguro en una transición
de costo cero. Luego,
\begin{align*}
    D_{\mathrm{DAgger}} &= \gamma S_{\mathrm{DAgger}}, \\
    \bigl(1-\gamma(1-\epsilon)-\epsilon\gamma^2\bigr)S_{\mathrm{DAgger}}
        &=\epsilon, \\
    J(\pi_{\mathrm{DAgger}})&=
        \boxed{\frac{\epsilon}{(1-\gamma)(1+\gamma\epsilon)}}.
\end{align*}

\subsection*{b) Comparación numérica}

Para $\epsilon=0.05$ y $\gamma=0.90$, se obtiene
\[
    J(\pi_{\mathrm{BC}})=\frac{0.05}{0.1(0.145)}\approx 3.448,
    \qquad
    J(\pi_{\mathrm{DAgger}})=\frac{0.05}{0.1(1.045)}\approx 0.478.
\]
El factor de reducción es
\[
    \frac{J(\pi_{\mathrm{BC}})}{J(\pi_{\mathrm{DAgger}})}
    =\frac{1.045}{0.145}\approx\boxed{7.21}.
\]
DAgger incurre en aproximadamente el 13,9\% del costo de BC,
lo que equivale a una reducción del 86,1\%.

\subsection*{c) Costo esperado en un horizonte finito}

Numérese del $1$ al $T$ cada transición desde el estado inicial seguro.
El experto permanece en el camino seguro y paga costo cero. Con BC,
las demostraciones solo cubren ese estado. Una caída lleva al robot a
un estado que el experto nunca visita desde el inicio y donde BC
continúa pagando. En el paso $n$ se incurre en costo $1$ si ocurrió
alguna caída en los primeros $n$ pasos, incluido el paso actual.
Evitar la caída durante los $n$ pasos tiene probabilidad $(1-\epsilon)^n$. Por
linealidad de la esperanza y la cota de la unión,
\begin{align*}
    J_T(\pi_{\mathrm{BC}})
      &= \sum_{n=1}^{T}\bigl[1-(1-\epsilon)^n\bigr] \\
      &\leq \sum_{n=1}^{T}\min\{1,n\epsilon\}
      \leq \min\left\{T,\frac{\epsilon T(T+1)}{2}\right\}.
\end{align*}
Para $\epsilon T\ll 1$,
$1-(1-\epsilon)^n=n\epsilon+O(n^2\epsilon^2)$, de donde
$J_T(\pi_{\mathrm{BC}})=\epsilon T(T+1)/2+O(\epsilon^2T^3)$.
El primer error puede afectar todos los pasos restantes, lo que
produce la dependencia cuadrática $O(\epsilon T^2)$. Para $\epsilon$
fijo y $T$ grande, el costo se acerca a $T$ en términos relativos,
porque $T-J_T(\pi_{\mathrm{BC}})\leq (1-\epsilon)/\epsilon$.

DAgger incorpora al conjunto de entrenamiento estados visitados por
la política en ejecución, etiquetados con acciones del experto
\cite[seccs.~2--4]{ross2011}. En este problema, aprende a salir de la
zanja en el siguiente paso, con costo cero. Sea $p_n$ la probabilidad
de estar en el camino seguro justo antes de la transición $n$.
Entonces $p_1=1$ y $p_{n+1}=1-\epsilon p_n$, pues se empieza el paso
siguiente en la zanja únicamente si se acaba de caer. Cada caída
cuesta una unidad, por lo que
\[
    J_T(\pi_{\mathrm{DAgger}})=\epsilon\sum_{n=1}^T p_n
         \leq \epsilon T.
\]
La cota es lineal porque cada error afecta el costo de una sola
transición.
